\section{Scope and Limitations}
\label{sec:discussion}

\subsection{Threat Model and Security Scope}

The target deployment assumes an organizational CA or identity provider with
enrolled devices and users; the backend trusts certificate-derived identity only
when the reverse proxy and service boundary are correctly configured. It handles
replayed field submissions through \textit{client\_event\_id} idempotency and
limits tablet reachability through local edge APIs and network segmentation.

Several threats remain outside the prototype evaluation. A stolen or compromised
tablet can be revoked only when nodes regain contact with the authority
infrastructure or receive an updated trust bundle, so instantaneous revocation
during total disconnection is not provided. A stolen vehicle node is more serious
because it may hold cached data, and if it is the command edge it also holds
operational authority. The architecture relies on device encryption, operational
fencing, and audit rather than Byzantine tolerance: because authority is explicit,
a captured or buggy command edge can poison the incident journal while it holds
authority, a central trade-off of making command authority explicit.

The design also does not address a compromised organizational CA, Byzantine
command behavior, physical tamper resistance beyond deployment controls, or
audit-log integrity against a fully compromised node. These are deployment and
hardening questions, not properties established by the emulation.


\subsection{Evaluation Boundaries and Non-claims}
\label{sec:non-claims}

The evidence for each invariant matches the appropriate verification method.
\textit{SingleAuthority} and \textit{SingleCommand} express properties over all
possible interleavings of concurrent authority claims---properties that model
checking establishes exhaustively over 38.56\,M distinct states but that a
service-path test can only sample. The prototype's service path is single-edge
because core-to-edge journal bootstrap is not implemented, so multi-edge single
authority is verified at the model level only; the two-edge Hypothesis handoff
model covers prefix bootstrap and stale-authority rejection, and semantic
rejection remains outside the prototype. \textit{EpochAuthorityCoupling},
\textit{NoForkedJournal}, \textit{DurabilityAcrossPromotion}, and
\textit{LocalIdempotency} have both model verification and live service-path
witnesses.

The prototype is scoped to measure protocol behavior under controlled partition
conditions. The emulation isolates the authority, outbox, and fencing mechanisms
from deployment variables---physical mesh transport, security overlay cost, and
Android client integration---that do not bear on the authority model's
correctness. Accordingly, it does not claim physical Rajant mesh behavior,
measured WireGuard or mTLS overhead, or end-to-end Android validation, and the
workload is single-incident and limited-scale. The authority model's safety
properties are established by the TLA+ verification; the prototype demonstrates
that the protocol can be implemented and operated under partition, and the
Hypothesis state-machine tests check that the implementation's behavior agrees
with the model on the sampled traces.

